\documentclass[11pt]{article}
\usepackage[a4paper,margin=1in]{geometry}
\usepackage{amsmath,amssymb}
\usepackage{hyperref}
\usepackage{graphicx}
\usepackage{booktabs}
\usepackage{microtype}

\title{The Compute-Allocation Frontier: A Living, Reproducible Survey of LLM Reasoning and Test-Time Compute}
\author{Pengyi Research\\ \small P-Research (living corpus pipeline) \\ \small \texttt{github.com/p-research}}
\date{Draft v0.1 -- generated \today}

\begin{document}
\maketitle

\begin{abstract}
Large language models are shifting from train-time to test-time scaling: reasoning models now spend deliberate compute at inference through long chains of thought, learned search, and verifier-guided refinement, trained via reinforcement learning with verifiable rewards (RLVR). Yet the field's existing surveys are static snapshots written by individual teams at a single moment, and they treat test-time compute in isolation from the rest of frontier AI. This paper presents a living, reproducible survey of LLM reasoning and test-time compute, generated by a continuously-updating arXiv corpus pipeline that sweeps six pillars weekly into an append-only, evidence-linked paper database. We organize the field under a unifying compute-allocation lens---where compute is spent across training versus inference, in the model versus the environment, and on representation versus reward. Every claim links to a verifiable paper record, making the survey auditable and updatable rather than a frozen artifact.
\end{abstract}

% ---- 01-introduction.md ----
\section*{\S1 Introduction}

\begin{quote}Draft v0.1. Claims cite Evergreen corpus records (\texttt{data/papers.jsonl}); only full-text-verified records may survive into the final version.\end{quote}

\section{1.1 From train-time to test-time compute}

For most of the deep-learning decade, scaling meant training: more data,

more parameters, more FLOPs at the optimization step. Since 2023 the field

has been re-negotiating that contract. Reasoning models now spend deliberate

compute *at inference*: long chains of thought, learned search, and

verifier-guided refinement, trained with reinforcement learning on

verifiable rewards (RLVR). The frontier has moved from "a bigger model" to

"a better allocation of compute" --- and the allocation question now spans

training vs. inference, model vs. environment, and representation vs. reward.


The shift is visible even in abstract-level signals. In our corpus of 503

frontier-AI papers (2023--2026, six pillars), RLVR/GRPO appears in 4

full-text-verified papers from 2024, 18 from 2025, and 6 in the first part

of 2026 --- while classic chain-of-thought language rises more slowly

(2 $\rightarrow$ 12 $\rightarrow$ 4). Test-time scaling and verifier/process-reward-model methods

follow the same trajectory. (Full corpus: \texttt{data/papers.jsonl}; verified

subset: 45 records, method-tag match rate 99\%.)


\section{1.2 Why a living survey}

Existing surveys on LLM reasoning and test-time compute --- e.g. the RLVR

survey \texttt{arXiv:2501.09686} and the test-time-compute survey \texttt{arXiv:2501.02497}

--- are static snapshots written by a single team at one moment. They are

excellent maps, but the territory moves weekly. A static survey cannot show

you that multi-agent methods doubled in the verified corpus between 2025

and 2026, or that a method you assumed was reasoning-specific is now

converging across pillars.


This survey is generated by a reproducible, continuously-updating corpus

pipeline (code on GitHub, append-only database, deterministic structuring,

auditable evidence chain from arXiv record to survey claim). It accumulates

evidence instead of rotting: every claim links to database records, and the

weekly digests archive the evolving frontier.


\section{1.3 The compute-allocation lens}

We organize the field under one question: \textbf{where is compute spent?}

Three axes organize the six pillars of our taxonomy:


1. \textbf{train-time vs. test-time} --- pre-training and post-training versus

inference-time reasoning budgets;

2. \textbf{in model vs. in environment} --- parameters and activations versus

tools, memory, and world interaction;

3. \textbf{on representation vs. on reward} --- learning what the world is like

versus learning what "good" means.


Under this lens, LLM reasoning (pillar 1), agentic deep research (pillar 2),

efficient training \& inference (pillar 3), RL/alignment/safety (pillar 4),

multimodal/world models (pillar 5), and quant$\times$AI (pillar 6) are not separate

fields but points on one tradeoff curve. \S5 develops this convergence with

corpus evidence.


\section{1.4 Contributions}

1. A \textbf{living corpus pipeline} (503 papers, 2023--2026; 45 lead-pillar

papers full-text verified at the time of writing) with deterministic

structuring and an auditable evidence chain.

2. A \textbf{compute-allocation taxonomy} that places test-time compute in the

context of the other five pillars.

3. \textbf{Empirical trend sections} derived from the corpus: method migration

(RLVR 4$\rightarrow$18 verified papers 2024$\rightarrow$2025), benchmark saturation (MATH/GSM8K

dominance), and cross-pillar convergence (\S5--\S6).

4. A public, continuously-updated artifact rather than a frozen PDF:

digests, database, and survey outline update weekly.


\section{1.5 Scope and honesty}

\begin{itemize}
\item Coverage: arXiv, six pillars, 2023 onward; English-language, abstract
\end{itemize}
level for unverified records, full text for verified records.

\begin{itemize}
\item Not covered: closed commercial systems without papers; non-ArXiv venues
\end{itemize}
beyond the verified references we cite.

\begin{itemize}
\item Every trend number in this draft is a *corpus* statement, not a claim
\end{itemize}
about the field's ground truth; where full-text verification is pending,

we say so.



% ---- 02-method.md ----
\section*{\S2 Method: A Reproducible, Continuously-Updating Corpus Pipeline}

\begin{quote}Draft v0.1 --- every claim must eventually cite \texttt{data/papers.jsonl} records. This section documents the pipeline that *generates* this survey.\end{quote}

\section{2.1 Design principles}

1. \textbf{Deterministic and auditable.} Every structuring step is a keyword/regex

rule, not an LLM call. A claim about a paper can be re-derived from its

abstract and full text by anyone.

2. \textbf{Append-only evidence.} The paper database (\texttt{data/papers.jsonl}) is

append-only; verification metadata is merged via atomic rewrites. The

evidence chain runs from arXiv entry $\rightarrow$ structured record $\rightarrow$ survey claim.

3. \textbf{Living, not frozen.} A weekly cron sweeps the frontier; historical

backfills cover 2023 onward. Sections cite the corpus *state*, and digests

archive the evolving frontier.

4. \textbf{Honest about signal strength.} Abstract-level tags are signals;

full-text verification is a separate gate. Only verified records may back

survey claims.


\section{2.2 Corpus construction}

\begin{itemize}
\item \textbf{Sources}: the arXiv API (export.arxiv.org), six pillar queries over
\end{itemize}
categories cs.AI, cs.LG, cs.CL, cs.CV, cs.CR, cs.MA, cs.DC, q-fin.

\begin{itemize}
\item \textbf{Windows}: weekly frontier sweep (last 21 days) + non-overlapping
\end{itemize}
historical backfill windows (0--360, 360--720, 720--1080, 1080--2160 days).

\begin{itemize}
\item \textbf{Rate limiting \& robustness}: $\geq$3s between requests, retries with backoff,
\end{itemize}
time-boxed cache, and salvage of complete \texttt{<entry>} blocks from truncated

responses.

\begin{itemize}
\item \textbf{Size}: 503 records (2023--2026) at the time of writing; 45 lead-pillar
\end{itemize}
records full-text verified.


\section{2.3 Deterministic structuring}

Each arXiv entry is mapped to a record with:


\begin{itemize}
\item \textbf{pillar} --- the query of origin (six-pillar taxonomy);
\item \textbf{methods} --- 22 keyword families (RLVR/GRPO, Verifier/PRM, CoT, MCTS,
\end{itemize}
MoE, distillation, quantization, KV cache, long context, preference

optimization, interpretability, safety/jailbreak, multi-agent, tool use,

memory/RAG, computer use, deep research, world models, video generation,

VLM, quant/trading); short acronyms use word-boundary matching to avoid

false positives (e.g. "cot" inside "scotland");

\begin{itemize}
\item \textbf{benchmarks} --- a curated list (AIME, GSM8K, MATH, GPQA, SWE-bench,
\end{itemize}
MMLU, ARC-AGI, AgentBench, GAIA, \ldots) with word-boundary rules for short

names;

\begin{itemize}
\item \textbf{models} --- regex patterns for common model families;
\item \textbf{key\_results} --- sentences carrying result markers (outperform, achieves,
\end{itemize}
state-of-the-art, \ldots).


\section{2.4 Full-text verification (core-corpus gate)}

\begin{itemize}
\item \textbf{Sources}: ar5iv (LaTeXML HTML) first; arXiv native HTML for papers ar5iv
\end{itemize}
has not converted. arXiv abstract pages masquerading as fulltext are

rejected by \texttt{ltx\_} marker detection.

\begin{itemize}
\item \textbf{Procedure}: pull full text $\rightarrow$ re-run the deterministic taggers on it $\rightarrow$
\end{itemize}
record \texttt{verified} plus matched-method overlap vs. the abstract-level tags.

\begin{itemize}
\item \textbf{Status}: 45 lead-pillar records verified; source split ar5iv/arxiv-html
\end{itemize}
recorded per record.


\section{2.5 Citation tracking (in progress)}

Semantic Scholar metadata (citationCount, influentialCitationCount) is

fetched per arXiv id and stored in \texttt{data/citations.jsonl}; anonymous access

degrades gracefully and records are only persisted on success.


\section{2.6 Reproducibility artifacts}

\begin{itemize}
\item \texttt{presearch} CLI (stdlib-only Python, $\geq$3.10): \texttt{weekly}, \texttt{backfill},
\end{itemize}
\texttt{verify}, \texttt{citations}, \texttt{db stats}, \texttt{survey}.

\begin{itemize}
\item GitHub Actions cron (weekly) commits new research back to the repository.
\item MIT code, CC BY 4.0 data.
\end{itemize}


% ---- 03-foundations.md ----
\section*{\S3 Foundations: From Chain-of-Thought to RLVR}

\begin{quote}Draft v0.1. \texttt{evg-*} ids cite Evergreen corpus records; only full-text-verified records are cited as evidence.\end{quote}

\section{3.1 Chain-of-thought as the seed}

Chain-of-thought (CoT) prompting showed that allocating inference tokens to

intermediate reasoning improves accuracy without touching weights. In our

verified corpus, CoT language appears in 2 verified 2024 papers and 12

verified 2025 papers --- but increasingly as a *substrate* of learned

reasoning rather than as a prompting trick: reasoning traces became the

object of optimization, not just an artifact of prompts.


\section{3.2 Verifiable rewards}

The RLVR recipe --- reinforcement learning where the reward is a

mechanically checkable function (a final answer, a unit test, a compiler

verdict) --- removes the need for a learned reward model on many tasks.

Verifier/process-reward-model methods grow in the verified corpus from 1

(2024) to 9 (2025) to 5 (partial 2026). Representative verified records:


\begin{itemize}
\item \texttt{evg-e04110c6e72bd0da} (arXiv:2408.15565, 2024): self-improving code-assisted
\end{itemize}
mathematical reasoning --- an early bridge between code execution and

verifiable rewards.

\begin{itemize}
\item \texttt{evg-a1327781cf590ac8} (arXiv:2508.16921, 2025): reward/verifier machinery
\end{itemize}
applied beyond math, to affective-alignment diagnosis --- a sign of the

stack spreading across pillars.


\section{3.3 GRPO, PPO, and the policy-optimization zoo}

Within RLVR, group-relative policy optimization (GRPO) became the

workhorse because it drops the critic model, cutting memory and compute ---

an *efficiency* choice that doubles as an *allocation* choice (reward vs.

representation). Our corpus shows preference-optimization language in 3

verified 2024 papers, 20 verified 2025 papers, and 7 in 2026 --- the fastest

growth of any method family in the verified subset. Policy optimization now

appears far outside its origin: in embodied robot manipulation

(\texttt{evg-4e48d42ef4f7ad7e}, BATON), in inverse RL for control (\texttt{evg-3f12f30b5bc710da}),

and in dialogue model-based RL (\texttt{evg-964669cd2625570d}).


\section{3.4 Search and self-correction}

Test-time search (MCTS-style tree search, beam variants, best-of-N) turns

inference into an optimization loop over candidate traces. Verified

examples: \texttt{evg-4e48d42ef4f7ad7e} (subtask-level exploration composed into

long-horizon plans), \texttt{evg-37a3465e623bd6ff} (AutoSR: symbolic regression by

searching "research states"), \texttt{evg-6742b060cb568d54} (Le Critique:

privileged value functions guiding LLM reinforcement learning), and

\texttt{evg-f945ded0c59d6837} (PuzzleJAX, a benchmark designed for

reasoning-and-learning under search). Search methods in the verified corpus

grow 1 (2024) $\rightarrow$ 6 (2025) $\rightarrow$ 3 (partial 2026).


\section{3.5 Reward hacking and the verifier-quality wall}

As verifiable rewards scale, so does gaming them. The safety pillar already

feeds back into reasoning: verified record \texttt{evg-204795ac0a186819} (What Do

Compliance Detectors Read?) audits activation probes and guard models ---

detector *representations* are now part of the reasoning stack's failure

modes. We flag an open problem here: **verifier quality is the new

bottleneck** --- a claim we develop with corpus evidence in \S7.


\section{3.6 What the corpus shows}

\begin{table}[h]\centering\begin{tabular}{llll}
\hline
Method family (verified full-text) & 2024 & 2025 & 2026 (partial) \\
\hline
RLVR / GRPO & 4 & 18 & 6 \\
Preference optimization & 3 & 20 & 7 \\
Chain-of-thought & 2 & 12 & 4 \\
Verifier / PRM & 1 & 9 & 5 \\
Search / MCTS & 1 & 6 & 3 \\
Test-time scaling & 0 & 7 & 5 \\
\hline\end{tabular}\end{table}


*Corpus caveat: weekly sweeps weight recent years; 2026 is a partial year.

These are corpus frequencies, not field-wide prevalence.*


\section{3.7 Open threads for \S4}

\begin{itemize}
\item How do verifiers and search interact with *inference-time scaling laws*?
\item Where does the budget go --- trace length, search width, or verifier depth?
\item What transfers from the reasoning pillar to agents, efficiency, and
\end{itemize}
quant$\times$AI (\S5)?



% ---- 04-test-time-compute.md ----
\section*{\S4 Test-Time Compute}

\begin{quote}Draft v0.1. \texttt{evg-*} ids cite Evergreen corpus records (full-text verified unless noted). Formal scaling-law results are summarized from the survey literature; corpus evidence covers the *system designs*.\end{quote}

\section{4.1 The allocation question}

Test-time compute is the budget a model spends at inference --- on tokens,

candidate traces, verifier calls, or tool interactions --- to raise answer

quality without touching weights. The design space decomposes into three

questions we use to organize this section:


1. \textbf{Where does search happen?} --- tree search, beam variants, best-of-N,

or learned routing over candidate traces.

2. \textbf{What judges a trace?} --- outcome verifiers, process reward models,

learned value functions, or external tools.

3. \textbf{How does the budget scale?} --- inference-time scaling laws and the

regime where more compute stops helping.


\section{4.2 Search: from fixed traces to composed plans}

Verified corpus examples show search moving from "sample many answers" to

"search over structured states":


\begin{itemize}
\item \texttt{evg-4e48d42ef4f7ad7e} (BATON) --- subtask-level exploration: each subtask is
\end{itemize}
explored in a cheap short-horizon regime, and long-horizon trajectories

are *composed* from stored solutions, turning multiplicative exploration

cost (T to the power K) into additive (T times K).

\begin{itemize}
\item \texttt{evg-37a3465e623bd6ff} (AutoSR) --- symbolic regression framed as searching
\end{itemize}
"research states", with a verifier-governed transition model.

\begin{itemize}
\item \texttt{evg-6742b060cb568d54} (Le Critique) --- privileged value functions steer LLM
\end{itemize}
policy search, reintroducing critic-shaped guidance into the RLVR loop.

\begin{itemize}
\item \texttt{evg-f469e87d0194dc45} (Learning from Diverse Reasoning Paths) --- routing and
\end{itemize}
collaboration across diverse reasoning paths rather than single-trace

decoding.

\begin{itemize}
\item \texttt{evg-f945ded0c59d6837} (PuzzleJAX) --- a benchmark built for reasoning *and
\end{itemize}
learning* under search, i.e. search as a first-class evaluation axis.


Reading: search is absorbing the *planning* role that classical AI assigned

to planners --- but now the search space is language traces and the value

function is a learned or verifiable judge (\S4.3).


\section{4.3 Verifiers and process reward models}

The judge is the new bottleneck (\S3.5). Verified examples:


\begin{itemize}
\item \texttt{evg-58ed686e9dc31903} (GRIP) --- grounded reasoning via
\end{itemize}
information-restricted premises: restricting what a trace may assume is

itself a verifier mechanism.

\begin{itemize}
\item \texttt{evg-e804bfb18640fc6e} (ClawGym II) --- black-box RL on an agent harness: the
\end{itemize}
environment as the ultimate verifier for agentic policies.

\begin{itemize}
\item \texttt{evg-204795ac0a186819} (Compliance Detectors) --- audits of activation probes
\end{itemize}
and guard models; detector *representation* quality becomes part of the

reasoning stack's failure modes.


For the formal taxonomy of outcome vs. process reward models and their

failure modes (reward hacking, credit assignment over long traces), the

reader should consult \texttt{arXiv:2501.09686} and \texttt{arXiv:2501.02497}; our corpus

contributes the empirical observation that verifier-style machinery now

appears in 15 verified lead-pillar records --- and in guard-model audits far

outside math reasoning.


\section{4.4 Inference-time scaling laws}

The quantitative literature (compute-optimal inference, budget

allocation, self-refinement limits) is summarized by the test-time-compute

survey \texttt{arXiv:2501.02497}; several more recent scaling-law surveys exist,

and we cite only ids verified against the live arXiv API.

Our corpus's contribution is *systemic* rather than formal: verified

records show test-time-scaling language in 7 papers from 2025 and 5 from

the partial year 2026 --- and, importantly, it has begun to span pillars

(lead: 3, efficiency: 1, quant: 1 at abstract level). The scaling-law

question we can now ask the corpus: **does the frontier treat test-time

compute as a lever (spend more when it helps) or a budget (allocate a fixed

envelope)?** --- a question flagged for \S7.


\section{4.5 Summary}

\begin{table}[h]\centering\begin{tabular}{lll}
\hline
Mechanism & Verified examples & What it allocates \\
\hline
Structured search & BATON, AutoSR, Le Critique & trace space $\times$ time \\
Verifier / PRM & GRIP, ClawGym II, Compliance Detectors & judge quality \\
Routing/collaboration & Diverse Reasoning Paths & width vs. depth \\
Benchmark pressure & PuzzleJAX & evaluation budget \\
\hline\end{tabular}\end{table}


*All records full-text verified; see \S2.4 for the verification procedure.*



% ---- 05-cross-pillar.md ----
\section*{\S5 Cross-Pillar Convergence under a Compute-Allocation Lens}

\begin{quote}Draft v0.1. Counts are abstract-level corpus frequencies (503 records); pillar assignment = query of origin. Verified subsets are noted where used.\end{quote}

\section{5.1 Method families that refuse to stay in one pillar}

A method family that appears in several pillars at once is either a fad, a

foundation, or a bridge. The corpus matrix (papers per pillar) for families

spanning $\geq$ 5 pillars:


\begin{table}[h]\centering\begin{tabular}{lllllll}
\hline
Family & Reasoning & Agents & Efficiency & RL/Align & Multimodal & Quant \\
\hline
RLVR / GRPO & 29 & 18 & 5 & 9 & 2 & 7 \\
Preference optimization & 14 & 1 & 8 & 24 & 17 & 4 \\
VLM & 7 & 3 & 8 & 4 & 57 & 2 \\
Interpretability & 8 & 4 & 5 & 31 & 2 & 7 \\
Memory / RAG & 9 & 8 & 16 & 2 & 5 & 9 \\
Deep Research & 4 & 8 & 1 & 1 & 3 & 7 \\
Safety / Jailbreak & 11 & 6 & 1 & 12 & 5 & 4 \\
Distillation & 3 & 1 & 37 & 1 & 5 & 3 \\
Video generation & 6 & 2 & 5 & 1 & 19 & 5 \\
\hline\end{tabular}\end{table}


*Reading guide: rows are allocation strategies; columns are where compute is

spent. The table is generated from \texttt{data/papers.jsonl} and regenerated on

every sweep.*


\section{5.2 RLVR is the lingua franca}

RLVR/GRPO is the strongest convergence signal in the corpus: present in all

six pillars, 29 + 18 + 9 + 7 + 5 + 2 = 70 records. The recipe --- optimize a

policy against a verifiable reward --- no longer belongs to math reasoning.

Embodied agents (BATON), black-box agent harnesses (ClawGym II), and

on-chain/finance pipelines all adopt it. Under the compute-allocation lens,

RLVR is a *reward-side* allocation: you spend compute deciding what "good"

means, and the policy follows.


\section{5.3 The efficiency pillar is the mirror image}

Where reasoning spends compute at inference, efficiency spends it at design

time --- distillation (37 records), quantization (27), MoE (12), KV-cache

work (7) --- to *buy back* the inference budget. The two pillars are not

opposed; they are the same curve. Distillation from a reasoning teacher to

a cheap student is literally test-time-compute crystallized into weights.

This is the sharpest single observation of the compute-allocation lens,

and \S6 quantifies it with trend data.


\section{5.4 Memory and verifiers meet in the middle}

Memory/RAG spans all six pillars (9/8/16/2/5/9) and verifier/PRM language

leaks from reasoning into efficiency and alignment. The mechanism is

shared: both are ways to *move compute out of the model* --- into stored

context (memory) or into an external judge (verifier). Agents, long-context

models, and retrieval-augmented pipelines are the same allocation move at

different timescales.


\section{5.5 Quant$\times$AI as a convergence sink}

Quant$\times$AI hosts 56 quant/trading records but also imports the frontier

stack: RLVR (7), deep research (7), interpretability (7), memory (9),

video generation (5), preference optimization (4). Finance is where the

reasoning pillar's *verifiable reward* (P\&L, risk limits, backtests) is

native --- which makes it a natural stress test for \S7's open problems.


\section{5.6 Caveats}

\begin{itemize}
\item Abstract-level tags overstate co-occurrence: a paper mentioning
\end{itemize}
"safety" is not a safety paper. Full-text verification (\S2.4) is the

gate that turns these frequencies into citable claims; verified-subset

versions of this table ship with the final draft.

\begin{itemize}
\item Pillar assignment follows the query of origin; cross-pillar numbers
\end{itemize}
therefore measure *query overlap* plus true method migration. We report

both effects rather than hiding either.


\section{5.7 What \S5 claims, once verified}

1. RLVR/GRPO is the most widely shared method family across pillars.

2. Efficiency methods are the counter-allocation to test-time compute.

3. Memory and verifiers are the two dominant "out-of-model" compute moves.

4. Quant$\times$AI is the frontier stack's most demanding verifiable-reward testbed.



% ---- 06-empirical-trends.md ----
\section*{\S6 Empirical Trends from the Living Corpus}

\begin{quote}Draft v0.1. All numbers are corpus statements, regenerated from \texttt{data/papers.jsonl} on every sweep. Verified subsets are preferred where available; abstract-level counts are labeled as such.\end{quote}

\section{6.1 Method migration: the RLVR wave (verified)}

Full-text-verified records show the sharpest single trend in the corpus:

RLVR/GRPO quadrupled from 4 (2024) to 18 (2025) verified papers, with 6

already in the partial year 2026; preference-optimization language grew

3 $\rightarrow$ 20 $\rightarrow$ 7. Chain-of-thought (2 $\rightarrow$ 12 $\rightarrow$ 4) grew more slowly, consistent

with the \S3 reading that CoT is being absorbed as a *substrate* of learned

reasoning rather than remaining a prompting technique.


Abstract-level counts across all six pillars tell the same story with more

noise (RLVR: 17/18/21/14 over 2023-2026) --- the 2023 elevation is a

backfill-window artifact (\S6.4).


\section{6.2 Benchmark concentration --- and where the hard ones hide}

In verified lead-pillar full texts, three benchmarks dominate: MATH (13),

DROP (12), GSM8K (8). High-difficulty suites (AIME: 2, MMLU: 2) appear far

less often --- but note the *location* asymmetry: AIME/GPQA almost never

appear in abstracts, while MATH/GSM8K do. Abstract-level benchmark mining

therefore systematically under-counts hard benchmarks; full-text

verification corrects it. For survey purposes this matters: a reader

judging "reasoning progress" from abstracts alone will see an easier field

than the papers actually evaluate on.


\section{6.3 The frontier moves weekly}

The weekly digests (\texttt{data/weekly/}) archive the moving frontier. The

2026-W34 sweep (the first) ingested 69 papers across the six pillars with

zero degraded queries; convergence signals that week included VLM spanning

4 pillars (12 papers) and preference optimization spanning 4 (10 papers).

Each digest carries the sweep's convergence signals and degraded-pillar

log, so trend claims in this survey can be audited week by week.


\section{6.4 Corpus-construction biases (reported, not hidden)}

1. \textbf{Window weighting.} Weekly sweeps cover the last 21 days; backfills

cover fixed windows (0-360, 360-720, 720-1080, 1080-2160 days). Recent

years are over-represented; 2026 is partial.

2. \textbf{Query-of-origin pillar assignment.} A paper's pillar follows the

query that found it, not a content classifier. Cross-pillar counts

measure query overlap plus true migration (\S5.6).

3. \textbf{Abstract-level noise.} Keyword tags overstate co-occurrence; only

full-text-verified records (45, method match rate 99\%) back hard claims.

4. \textbf{arXiv-only.} Closed-model reports, non-English venues, and

non-arXiv publications are out of scope (\S1.5).


\section{6.5 What \S6 contributes}

\begin{itemize}
\item A *method-migration* estimate with verified backing (RLVR 4$\rightarrow$18).
\item A *benchmark-visibility asymmetry* (hard benchmarks hide in full text).
\item An *auditable cadence* --- every trend links to weekly digests and DB
\end{itemize}
records, so the survey's numbers can be re-derived, not just re-cited.



% ---- 07-open-problems.md ----
\section*{\S7 Open Problems, Risks, and Outlook}

\begin{quote}Draft v0.1. Open problems are ordered by how directly the corpus can test them.\end{quote}

\section{7.1 Verifier quality is the new bottleneck}

RLVR scales only as far as the reward signal is trustworthy. The corpus

already shows the failure mode migrating: verifier/guard-model audits

(\texttt{evg-204795ac0a186819}), affective-hallucination diagnosis

(\texttt{evg-a1327781cf590ac8}), and black-box agent-harness RL (\texttt{evg-e804bfb18640fc6e})

all circle the same question --- \textbf{how do you verify the verifier?}

Candidate research directions: adversarial verifier evaluation, reward

calibration under distribution shift, and process-reward attribution over

long traces.


\section{7.2 Test-time compute: lever or budget?}

\S4.4 asked whether the frontier treats test-time compute as a lever (spend

more when it helps) or a budget (fixed envelope). The corpus cannot answer

it yet --- scaling-law papers are under-represented in our queries, and the

formal literature lives in the surveys we cite (\texttt{arXiv:2501.02497}). A

corpus upgrade (dedicated scaling-law query family) is planned.


\section{7.3 The closed frontier is invisible to us}

Closed commercial systems (o-series, Claude, Gemini reasoning modes)

publish no full text; our corpus sees them only through third-party

evaluations. Any survey claim about "the frontier" is therefore

conditional on the open literature. Mitigation: position claims as

"open-literature frontier" and track closed-model evaluations as a

separate, clearly-labeled source class.


\section{7.4 Citation lag and novelty scoring}

2026 papers have immature citation graphs; Semantic Scholar coverage lags

submission by weeks. Novelty scoring from citations (\S2.5) must therefore

blend citation counts with *method-overlap novelty* (how unusual a

paper's method vector is within its pillar cohort) --- a direction we are

implementing.


\section{7.5 Reproducibility of a living artifact}

A "living survey" is only as living as its hosting. Risks: repo

abandonment, API changes, DB corruption. Mitigations in progress: full

DB checksums in digests, deterministic regeneration from raw arXiv

responses (cached), and periodic Zenodo snapshots of \texttt{data/}.


\section{7.6 Publication risk (arXiv endorsement)}

First-time submissions to cs.AI/cs.LG/cs.CL require endorsement, and arXiv

tightened its policy in January 2026. Status and fallback venues are

tracked in \texttt{positioning.md}; the survey's substance is designed to remain

valuable regardless of venue (open data + code + digests).


\section{7.7 Outlook}

Three bets close the survey:


1. \textbf{Reward-side compute} (RLVR + verifiers) keeps diffusing into every

pillar --- Quant$\times$AI is its most demanding testbed because P\&L is the

harshest verifier.

2. \textbf{Efficiency methods become reasoning methods' mirror}: distillation

crystallizes test-time compute into weights; the two pillars converge

on one allocation curve.

3. \textbf{Memory and verifiers merge}: the next reasoning stack will treat

stored context and external judges as one out-of-model compute

substrate.


The corpus will re-test these bets every Monday.




\section*{Acknowledgments}
This survey is generated by an open-source, continuously-updating arXiv
corpus pipeline. The full database, weekly digests, and verification
metadata are released under CC BY 4.0. All claims cite corpus records;
only full-text-verified records are treated as evidence.

\end{document}
